\documentclass[14pt]{beamer}
\usepackage{pepslides}
\title{Constitutive theory}
\lecturenumber{14}
\hypersetup{pdftitle={Lecture 14: Constitutive theory}}

\begin{document}

\titleslide

\begin{frame}[t]
\heading{From last time\ldots}
Using the concrete form for the Lagrangian density, the equations of motion are:
\begin{eqn}
\rho\frac{\partial v_{i}}{\partial t}-\frac{\partial T_{ij}}{\partial x_{j}}=0,
\end{eqn}
along with the traction-free boundary conditions:
\begin{eqn}
T_{ij}\hat{n}_{j}=0,
\end{eqn}
where we have defined the \textbf{first Piola--Kirchhoff stress} by:
\begin{eqn}
T_{ij}(\mathbf{x},t)=\frac{\partial W}{\partial F_{ij}}[\mathbf{x},\mathbf{F}(\mathbf{x},t)].
\end{eqn}
\end{frame}

\sectionslide{Material frame indifference}

\begin{frame}
For a simple hyperelastic material, the general form for the potential energy is:
\begin{eqn}
\mathcal{V}(t)=\int_{M}W[\mathbf{x},t,\bphi(\mathbf{x},t),\mathbf{F}(\mathbf{x},t)]\dd^{3}\mathbf{x}.
\end{eqn}
Using invariance under time-translation and superimposed rigid body motions, we showed that:
\begin{eqn}
W[\mathbf{x},\mathbf{F}(\mathbf{x},t)]=W[\mathbf{x},\mathbf{Q}(t)\mathbf{F}(\mathbf{x},t)].
\end{eqn}
How do we determine suitable classes of functions for the potential energy density?
\end{frame}

\begin{frame}
The \textbf{polar decomposition theorem} shows that
\begin{eqn}
\mathbf{F}=\mathbf{R}\mathbf{U}=\mathbf{V}\mathbf{R},
\end{eqn}
with $\mathbf{R}$ a rotation matrix, and $\mathbf{U}$ and $\mathbf{V}$ symmetric and positive-definite.

{\centering\includegraphics[height=0.56\paperheight]{L2F1.pdf}\par}

Proof in the handout, but \textbf{non-examinable}.
\end{frame}

\begin{frame}
Using the decomposition along with material frame indifference, we obtain:
\begin{eqn}
W(\mathbf{x},\mathbf{F})=W(\mathbf{x},\mathbf{R}\mathbf{U})=W(\mathbf{x},\mathbf{U}).
\end{eqn}
The potential energy density depends only on the stretching part of the deformation gradient.

Due to this result, we can always write:
\begin{eqn}
W(\mathbf{x},\mathbf{F})=U(\mathbf{x},\mathbf{C}),
\end{eqn}
where we have introduced the \textbf{right Cauchy--Green deformation tensor}:
\begin{eqn}
\mathbf{C}=\mathbf{F}^{T}\mathbf{F}.
\end{eqn}
This is a local measure of change of shape/volume, and from it we define the \textbf{finite strain tensor}:
\begin{eqn}
\mathbf{E}=\frac{1}{2}(\mathbf{C}-\mathbf{I}).
\end{eqn}
From here on, we call the potential energy density the \textbf{strain energy function} for the material.
\end{frame}

\begin{frame}
From the definition of the first Piola--Kirchhoff stress tensor, we have:
\begin{eqn}
T_{ij}=\frac{\partial W}{\partial F_{ij}}=\frac{\partial U}{\partial C_{kl}}\frac{\partial C_{kl}}{\partial F_{ij}}.
\end{eqn}
This implies that:
\begin{eqn}
T_{ij}=F_{ik}S_{kj},
\end{eqn}
where we have defined the \textbf{second Piola--Kirchhoff stress tensor}:
\begin{eqn}
S_{ij}=2\frac{\partial U}{\partial C_{ij}}.
\end{eqn}
This is a \textbf{symmetric tensor}, with this symmetry being equivalent to the \textbf{conservation of angular momentum} --
a fact you will work through in the first problem set.
\end{frame}

\sectionslide{Equilibrium and linearisation}

\begin{frame}
At equilibrium, the motion is time-independent. The equations of motion reduce to:
\begin{eqn}
\begin{gathered}
\frac{\partial T_{ij}^{0}}{\partial x_{j}}=0,\\[\eqgap]
T_{ij}^{0}\hat{n}_{j}=0.
\end{gathered}
\end{eqn}
We can choose the reference body to be the equilibrium configuration so that the \textbf{motion is the identity mapping}.

A solution of the problem is then given by a \textbf{stress-free equilibrium}.

We will assume such an equilibrium state for the first part of this course. But in the Earth, due to gravitation and
rotation, the equilibrium state is not stress-free.
\end{frame}

\begin{frame}
We consider the disturbance of a body about a stress-free equilibrium state:
\begin{eqn}
\begin{aligned}
\varphi_{i}(\mathbf{x},0)&=x_{i}+su_{i}^{0}(\mathbf{x}),\\[\eqgap]
v_{i}(\mathbf{x},0)&=sv_{i}^{0}(\mathbf{x}).
\end{aligned}
\end{eqn}
Here $s$ is a dimensionless perturbation parameter. We expand the motion in a Taylor series:
\begin{eqn}
\varphi_{i}(\mathbf{x},t)=x_{i}+su_{i}(\mathbf{x},t)+O(s^{2}),
\end{eqn}
where $\mathbf{u}$ is known as the \textbf{displacement vector}.

All we have to do is put this into the equations of motion, expand to first order in $s$ and see what happens!
\end{frame}

\begin{frame}
In linearising the stress tensor, we find:
\begin{eqn}
T_{ij}(\mathbf{x},t)=s\frac{\partial^{2}W}{\partial F_{ij}\partial F_{kl}}(\mathbf{x},\mathbf{I})\frac{\partial u_{k}}{\partial x_{l}}+O(s^{2}).
\end{eqn}
Motivated by this we introduce the \textbf{elastic tensor}:
\begin{eqn}
A_{ijkl}(\mathbf{x})=\frac{\partial^{2}W}{\partial F_{ij}\partial F_{kl}}(\mathbf{x},\mathbf{I}).
\end{eqn}
Due to the equality of mixed partial derivatives, we note the \textbf{hyperelastic symmetry}:
\begin{eqn}
A_{ijkl}=A_{klij}.
\end{eqn}
\end{frame}

\begin{frame}
Due to the stress-free equilibrium assumption, you will show that:
\begin{eqn}
A_{ijkl}=A_{jikl}=A_{ijlk}.
\end{eqn}
Along with the hyperelastic symmetry, this means there are \textbf{21 independent elastic moduli}.

The linearised stress depends only on the symmetrised displacement gradient:
\begin{eqn}
A_{ijkl}\frac{\partial u_{k}}{\partial x_{l}}=A_{ijkl}\frac{1}{2}\left(\frac{\partial u_{k}}{\partial x_{l}}+\frac{\partial u_{l}}{\partial x_{k}}\right).
\end{eqn}
Here we see the \textbf{linearised strain tensor}, which is the first-order term in the expansion of the finite strain.
\end{frame}

\begin{frame}
The final linearised equations of motion are:
\begin{eqn}
\rho\frac{\partial^{2}u_{i}}{\partial t^{2}}-\frac{\partial}{\partial x_{j}}\left(A_{ijkl}\frac{\partial u_{k}}{\partial x_{l}}\right)=0,
\end{eqn}
along with boundary conditions:
\begin{eqn}
\hat{n}_{j}A_{ijkl}\frac{\partial u_{k}}{\partial x_{l}}=0.
\end{eqn}
Note that the boundary conditions are applied on the reference boundary without need for further approximation.

We will see later that this is an example of a \textbf{wave equation}.
\end{frame}

\sectionslide{Material symmetries}

\begin{frame}
In general, the following relation, with $\mathbf{Q}$ a rotation matrix, does not hold:
\begin{eqn}
W(\mathbf{x},\mathbf{F}\mathbf{Q})=W(\mathbf{x},\mathbf{F}).
\end{eqn}
If it does, we say that $\mathbf{Q}$ is in the \textbf{material symmetry group} of $W$ at the point $\mathbf{x}$.

Material symmetries reflect the invariance of the strain energy with respect to the \textbf{orientation of stretching}.

On a crystal scale, they reflect lattice structures.
\end{frame}

\begin{frame}
If the material symmetry group is \textbf{all rotation matrices}, the material is \textbf{isotropic} at that point.

Recalling frame indifference, this is equivalent to:
\begin{eqn}
U(\mathbf{C})=U(\mathbf{Q}^{T}\mathbf{C}\mathbf{Q}).
\end{eqn}
But if $\mathbf{Q}$ is an arbitrary rotation matrix, we can select it so that the argument on the right-hand side is diagonalised.

Thus, in an isotropic material the strain energy depends on $\mathbf{C}$ through its eigenvalues.

From this, and some tedious algebra, we can show that:
\begin{eqn}
A_{ijkl}=\lambda\delta_{ij}\delta_{kl}+\mu(\delta_{ik}\delta_{jl}+\delta_{il}\delta_{jk}).
\end{eqn}
An isotropic material has only \textbf{two elastic moduli}, the \textbf{Lamé parameters} $\lambda$ and $\mu$!
\end{frame}

\begin{frame}
In an isotropic material with spatially homogeneous properties, the linearised equations about a stress-free
equilibrium state reduce to the form:
\begin{eqn}
\rho\frac{\partial^{2}\mathbf{u}}{\partial t^{2}}-(\lambda+\mu)\nabla(\nabla\cdot\mathbf{u})-\mu\nabla^{2}\mathbf{u}=\mathbf{0}.
\end{eqn}
This equation is seen in Part IB Classical Dynamics. We have now derived it from first principles, and in
this manner we can also see that it is \textbf{not suitable} for understanding seismic waves in the Earth.
\end{frame}

\begin{frame}
A material that is not isotropic is said to be \textbf{anisotropic}.

A special case are \textbf{transversely isotropic} materials, whose symmetry group comprises rotations about a fixed
axis $\hat{\bm{\nu}}$. In this case, it can be shown that the elastic tensor takes the form:
\begin{eqn}
\begin{aligned}
A_{ijkl}={}&(A-2N)\delta_{ij}\delta_{kl}+N(\delta_{ik}\delta_{jl}+\delta_{il}\delta_{jk})\\[\eqgap]
&+(F-A+2N)(\delta_{ij}\hat{\nu}_{k}\hat{\nu}_{l}+\hat{\nu}_{i}\hat{\nu}_{j}\delta_{kl})\\[\eqgap]
&+(L-N)(\hat{\nu}_{i}\hat{\nu}_{k}\delta_{jl}+\hat{\nu}_{j}\hat{\nu}_{k}\delta_{il}+\hat{\nu}_{j}\hat{\nu}_{l}\delta_{ik}+\hat{\nu}_{i}\hat{\nu}_{l}\delta_{jk})\\[\eqgap]
&+(A+C-2F-4L)\hat{\nu}_{i}\hat{\nu}_{j}\hat{\nu}_{k}\hat{\nu}_{l}.
\end{aligned}
\end{eqn}
\end{frame}

\begin{frame}
Here there are \textbf{five elastic moduli}, $A$, $C$, $F$, $L$ and $N$, in the notation introduced by Love.

With the symmetry axis along $\hat{\mathbf{x}}_{3}$, they are the components:
\begin{eqn}
A=A_{1111},\quad C=A_{3333},\quad F=A_{1133},\quad L=A_{1313},\quad N=A_{1212}.
\end{eqn}
In the isotropic case, $A=C=\lambda+2\mu$, $F=\lambda$ and $L=N=\mu$.

You do not need to learn the expression for this elastic tensor. And you will always be provided with the necessary form of an elastic
tensor within any exam question.
\end{frame}

\begin{frame}
An elastic material is a \textbf{solid} if the symmetry group is a sub-group of rotations.

It is a \textbf{fluid} if it comprises all volume-preserving linear transformations, so that:
\begin{eqn}
W(\mathbf{x},\mathbf{F})=V(\mathbf{x},J).
\end{eqn}
Physically, a solid's potential energy is affected by changes of shape and volume relative to a reference state,
while a fluid cares only about changes in volume.

The elastic tensor then reduces to:
\begin{eqn}
A_{ijkl}=\kappa\delta_{ij}\delta_{kl},
\end{eqn}
where $\kappa$ is the \textbf{bulk modulus} of the fluid.

Within the linearised theory, an elastic fluid is just an elastic solid with \textbf{vanishing shear modulus}.
\end{frame}

\begin{frame}[t]
\heading{What you need to know and be able to do}
\begin{enumerate}
\item[(i)] The statement of the polar decomposition theorem, and how this can be used to restrict the allowable form
of the strain energy function. The proof of the polar decomposition theorem is non-examinable.
\item[(ii)] What a stress-free equilibrium configuration is, and how to derive the linearised equations of motion for
small disturbances about such a configuration. In this course, you will become very familiar with these linearised
equations of motion, and it will be necessary to memorise their form (in the non-gravitating case discussed in this
lecture) for the exam.
\item[(iii)] The definition of the material symmetry group, and its physical interpretation.
\item[(iv)] The meaning of the terms elastic solid, elastic fluid, isotropic, anisotropic, and transversely isotropic.
You do not need to remember the form of the elastic tensor in any particular case, but you should memorise the
hyperelastic symmetry $A_{ijkl} = A_{klij}$ due to its great physical importance.
\end{enumerate}
\end{frame}

\end{document}
